\par\Needspace{6\baselineskip}\medskip
\hypertarget{native-110-article081}{}\pdfbookmark[2]{81.}{article-081}
\begin{center}81.\end{center}
If the eccentricity $e$ is given, the angle $\Pi$ is found from the equation
\[
\cos(A-\Pi)=-\frac{\cos(A-\tfrac12N-\tfrac12N')}{e\cos\tfrac12(N'-N)}
\]
once the auxiliary angle $A$ satisfying the equation
\[
\operatorname{tang}(A-\tfrac12N-\tfrac12N')=\frac{r'+r}{r'-r}\operatorname{tang}\tfrac12(N'-N)
\]
has been determined. The ambiguity that remains in determining the angle $A-\Pi$ from its cosine is inherent in the problem, which may therefore have two distinct solutions. Which to adopt and which to reject must be decided from other information. For this purpose at least an approximate value of \SourceLeaf{111}$\Pi$ must already be known. — Once $\Pi$ has been found, $p$ is calculated either by the formulas
\[
p=r(1+e\cos(N-\Pi))=r'(1+e\cos(N'-\Pi))
\]
or by the formula
\[
p=\frac{2rr'e\sin\tfrac12(N'-N)\sin(\tfrac12N+\tfrac12N'-\Pi)}{r'-r}\mathpunct{\cdot}
\]
